\section{Discussion}
\label{sec:discussion}

\subsection{Infrastructure as the contribution}

The methods in Table~\ref{tab:analyses} are mostly not new. What is new is that
they sit on shared components: one network type, one set of topology moves, one
parameter optimiser, one I/O layer, one set of chain diagnostics. A developer
adding a criterion writes the criterion.

Whether this lowers the barrier enough to change how network methods get built is
not something this paper can settle. The registry is exercised by seven methods,
all of them ours, and an interface that suffices for its author is weaker
evidence than one an outside contributor found sufficient. The design argument is
stated in Section~\ref{sec:architecture}; the measurement that would test it does
not exist yet.

Two design decisions seem worth keeping regardless. Making the validity of a
(data, model, criterion) triple a runtime query, rather than prose in
documentation, means the method set cannot drift from its description --
Table~\ref{tab:analyses} is generated from the registry. And separating the three
failure modes gives a user an actionable answer: change the data, wait for a
release, or reconsider the question.

\subsection{What the Python implementation does and does not buy}

Being a Python library makes some things straightforward that were awkward
before. An inferred network is an object that goes directly into a distance
function, a plotting call, or a pandas frame. A sweep over reticulation counts is
a loop. Simulation and inference compose without an intermediate file, which is
what makes the recovery check in Section~\ref{sec:examples} two lines.

It does not, by itself, make anything faster. \phynetpy's runtime advantage in
Section~\ref{sec:headtohead} comes from compiled kernels and from algorithmic
choices such as incremental rescoring, not from the choice of language -- and the
component measurements show the trade is not uniform: the graph core is faster on
incidence queries and slower on whole-graph traversal, using about twice the
memory. Claims that one language is faster than another are not what this section
is for.

Integration with the numerical and machine-learning ecosystem is often cited as a
reason to work in Python, and the architecture does make it practical: the
network is an ordinary object, likelihoods are ordinary callables, and a scorer
backed by PyTorch or JAX could be registered as a cell without changing anything
else. \phynetpy\ 1.0.0 contains no machine-learning functionality, no
differentiable likelihood and no GPU support beyond an optional CuPy path in the
marker likelihood. The integration is available rather than done.

\subsection{Reproducibility}

Two properties of the release matter for reproducing results. The version is
single-sourced and CI verifies that a clean install of both the wheel and the
source distribution reports it, so a paper citing a version is citing something
checkable. And the parts of an analysis that affect results -- the model
parameters, the criterion, the search preset, the seed -- are arguments at the
call site rather than defaults buried in a driver file.

The benchmark scripts for this paper follow the same discipline: each writes a
JSON file recording package versions, interpreter, platform, replicate counts and
seeds, and the tables and figures are generated from those files. This is why the
accuracy result in Section~\ref{sec:headtohead} is reported as it is. It emerged
from re-running a benchmark rather than from reading code, and a pipeline that
only regenerated favourable numbers would not have surfaced it.

\subsection{Limitations}

Relative to PhyloNet, \phynetpy\ 1.0.0 lacks maximum parsimony from gene trees
under diploid hybridization, and its marker maximum likelihood can score a
network but not search. Both are visible in \code{validity_matrix()}.

The accuracy deficit in Section~\ref{sec:headtohead} is the most consequential
open item. Across 30 trials of matched pseudo-likelihood inference PhyloNet
recovered the generating topology 16 times and \phynetpy\ 8, and the
\code{phylonet} preset intended to make the comparison like-for-like did not
close the gap. Anyone whose work depends on pseudo-likelihood network inference
should validate against PhyloNet on data with a known answer before relying on
\phynetpy\ for it.

Other limits are narrower. Across-site rate heterogeneity is not implemented, so
sequence analyses cannot fit a discretised gamma over sites. Simulation supports
only the \code{MSC} model; the allopolyploid model has no simulator. The network
simulator grafts reticulations onto a separately drawn tree rather than running a
birth--death--hybridization process, so it should not be used to study the
distribution of network shapes. Visualisation is ASCII text. There is no network
classifier for the graph-theoretic classes beyond level, and no summarisation of
\emph{sets} of networks, which is what a Bayesian run most wants. One compiled
kernel, \code{seq_engine_cy}, is built but not yet called. Installing from source
needs a C compiler. And every measurement in this paper comes from one machine.

\subsection{Future work}

The immediate items are the two unbuilt cells and a diagnosis of the accuracy
result. Beyond those: connecting the compiled sequence kernel; rate
heterogeneity; a birth--death--hybridization simulator; network classification
and set-level summarisation; and visualisation, where writing Tracer and NEXUS
output has so far been a substitute for having any of our own.

The architecture suggests directions we have not taken. A differentiable
pseudo-likelihood would make gradient-based optimisation over continuous
parameters available; a learned proposal distribution would fit the
\code{ProposalKernel} interface; approximate Bayesian computation would compose
with \code{simulate} directly. These are plausible because the pieces are
separable, not because any of them is underway.
